%! Author = Saeid Yazdani
%! Date = 6/22/2024


\section{\acs{DPS} Control Loop}
\label{sub:dps-control-loop}

All \ac{ATE} power supplies are programmable, and the output driver is inside a
control loop that will sense the target voltage at the \ac{DUT} site.
\emph{Remote Sense} is mandatory in the era of high supply currents and
low supply voltages – without remote sensing capabilities, there would be a
large unregulated voltage drop across the load board that would make the test
meaningless.

Systems such as \acp{ATE} which exhibit
significant wiring-related parasitics and variable \ac{DUT} load conditions must
use remote sensing schemes to compensate for voltage drops in the
cables between the power supply and the load, ensuring the correct voltage is
delivered at the load.

A comparison of control loops with and without remote sensing and the
associated errors is presented in
\autoref{fig:distributed-arrangement-ate-supply}:

\begin{itemize}[topsep=8pt,itemsep=4pt,partopsep=4pt, parsep=4pt]
    \item \textbf{Local Sense:}
    Without \emph{Remote Sense} and only relying on \emph{Local Sense}, we
    stabilize the supply voltage at the output of the source, and all wiring
    effects between source and \ac{DUT} are neglected.
    Such an arrangement is not suitable for precision measurements
    (\autoref{fig:distributed-arrangement-ate-supply}-A).

    \item \textbf{Basic Remote Sense:}
    If we can neglect all inductances of the sense lines because of very low
    settling times, we may succeed with a simple resistive approach
    (\autoref{fig:distributed-arrangement-ate-supply}-B).

    \item \textbf{Complex Remote Sense:}
    Because test times should be short, we need fast settling, and inductive
    components of the load board, socket and  \ac{DUT} can no longer be
    neglected (\autoref{fig:distributed-arrangement-ate-supply}-C).

\end{itemize}

\begin{figure}[htbp]
    \centering
    \printfitgraphics[0.62\textheight]{figures/villach-supply-analysis/distributed_arangement_ate_dut_supply}
    \caption[Comparison of Sensing Arrangements in Supplies]{
        The distributed arrangement of ATE source and DUT supply pin.
    }
    \label{fig:distributed-arrangement-ate-supply}
\end{figure}

While power supplies with remote sensing compensate for the wiring parasitic,
there are multiple challenges associated with such structures.
Implementing remote sensing requires additional wiring and components,
which will add complexity to the system design and in \ac{ATE} environments, and
requires dedicated wires on the load board to bring sense lines near the
\ac{DUT} socket.

This additional wiring introduces additional
susceptibility to noise and interference as incorrect or poor implementation of
remote sensing on any section of the \ac{ATE} to \ac{DUT} path can lead to
instability in the control loop and can potentially contribute to oscillations
and sporadic glitches seen in \autoref{ch:casestudy-from-ate-to-dut}.

Reducing analog components count and overall complexity of control loop in
favour of digital solutions can lead to a faster and more accurate transient
response.




